% TOP_PROBLEM: {"id": 30006877, "title": "Arithmeticity of Lattices in PU(n,1) for n at Least 4", "category_id": 17, "category": {"id": 17, "name": "group_theory", "display_name": "Group Theory", "description": "Problems about groups, group actions, representations, and related algebraic structures.", "slug": "group-theory", "order_index": 17, "created_at": "2026-07-31T15:26:25.671Z"}, "status": "open"}
% FIRST_POSTED: 2026-09-25
% ENTRY_KIND: statement_only
% !TeX program = lualatex
% Definition and sources only; this file is not an LLM solution attempt.
\documentclass[11pt]{article}
\usepackage[margin=25mm]{geometry}
\usepackage{fontspec}
\setmainfont{Latin Modern Roman}
\usepackage{amsmath,amssymb,mathtools}
\usepackage{unicode-math}
\setmathfont{Latin Modern Math}
\usepackage{xurl,hyperref}
\hypersetup{colorlinks=true,urlcolor=blue}
\setcounter{secnumdepth}{0}
\setlength{\parindent}{0pt}
\setlength{\parskip}{5pt}
\setlength{\emergencystretch}{3em}
\begin{document}
\section*{\texorpdfstring{Arithmeticity of Lattices in PU(n,1) for n at Least 4}{Arithmeticity of Lattices in PU(n,1) for n at Least 4}}\label{30006877}
\subsection{Definitions and mathematical statement}
Complex hyperbolic $n$-space is the space of negative complex lines for a Hermitian form of signature $(n,1)$; its holomorphic isometry group is $\operatorname{PU}(n,1)$. A lattice $\Gamma$ is discrete and has finite Haar covolume. Two subgroups are commensurable if their intersection has finite index in both. In the arithmeticity convention here, $\Gamma$ is commensurable, after conjugation, with the projection of an arithmetic subgroup of a ${\mathbb{Q}}$-algebraic group through a real-group homomorphism with compact kernel; finite central lifts to $\operatorname{SU}(n,1)$ give the equivalent formulation. The assertion is
\[
 n\ge4,\quad\Gamma<\operatorname{PU}(n,1)\text{ a lattice}
 \quad\Longrightarrow\quad\Gamma\text{ arithmetic}.
\]
Here $n$ is complex dimension. Noncompact finite-volume quotients, cocompact lattices, and torsion are all allowed, without restricting the arithmetic construction to a particular type of Hermitian form.
\par\smallskip\textit{Formulation and concept references:} \href{https://www.maths.dur.ac.uk/users/j.r.parker/img/Lattices.pdf}{[S1]}.

\subsection{Short English statement}
Is every finite-covolume lattice in the holomorphic isometry group of complex hyperbolic space arithmetic once the complex dimension is at least four?
\subsection{Sources}
\begin{itemize}
\item[S1]Complex hyperbolic lattices; undated author-hosted survey snapshot, not assigned a publication date from PDF metadata.
\url{https://www.maths.dur.ac.uk/users/j.r.parker/img/Lattices.pdf}.
\textit{Formulation source.}
\end{itemize}
\end{document}
